\exposetitle{XIII}{Regular elements of algebraic groups and of Lie algebras}
\label{XIII}
\begin{center}by A. Grothendieck\end{center}
{\renewcommand{\thefootnote}{0}\footnotetext{xy version of 1/12/08}}
\setcounter{footnote}{0}

\sgasection{1}{An auxiliary lemma on varieties with operators}
Let $S$ be a prescheme, $G$ an $S$-group prescheme acting on the left on an
$S$-prescheme $V$, $W$ a closed sub-$S$-prescheme of $V$, and $N$ its stabilizer
in $G$, the subgroup of $G$ whose points with values in an $S'$ over $S$ are the
$g\in G(S')$ such that $g\cdot W_{S'}=W_{S'}$. One equips $(\Sch/S)$ with the
faithfully flat quasi-compact topology, and identifies $G,V,M$ with the
corresponding sheaves (cf. IV).
% typo?: kept as printed: M, not introduced here, occurs in the source.
One shall therefore reason in the category of sheaves on $(\Sch)_{/S}$, and in
this section the expression ``locally'' refers to the topology just specified on
$(\Sch)_{/S}$. Note that $N$ is a sheaf; consider the quotient sheaf $G/N$.
One sees immediately that it is isomorphic to the following functor: to each
$S'$ over $S$, one associates the set of subsheaves $W'$ of $V_{S'}$ which are
locally conjugate to $W_{S'}$ under the group $G$. Let $X$ be the subsheaf of
$G/N\times_S V$ whose value, for every $S'$ over $S$, is the set of $(W',v)$,
where $W'$ is as above and $v$ is a section of $W'$ over $S'$ (thus a section of
$V_{S'}$ over $S'$). Let $Z$ be the inverse image of $X$ in $G\times_S V$, so
that one has the cartesian diagram
\[
\xymatrix{
 Z\ar[r]\ar@{^{(}->}[d]&X\ar@{^{(}->}[d]^i\\
 G\times_S V\ar[r]&G/N\times_S V,
}\tag{x}
\]
where $i$ is the canonical immersion, and the second horizontal arrow comes
from the canonical morphism $G\to G/N$. Since this associates to the point
$g\in G(S')$ the subsheaf $g\cdot W_{S'}$ of $V_{S'}$, one sees that $Z(S')$ is
the set of pairs $(g,v)\in G(S')\times V(S')$ such that
$v\in g\cdot W_{S'}(S')$. Consequently, $Z$ is isomorphic to the sheaf
$G\times_S W$, by means of the isomorphism
\[
 G\underset{S}{\times}W\isomto Z
\]
defined by $(g,w)\mto(g,g\cdot w)$. Thus the preceding cartesian diagram gives
the cartesian diagram
\[
\xymatrix{
 G\times_S W\ar[r]^q\ar[d]_\lambda&X\ar[d]^i\\
 G\times_S V\ar[r]&G/N\times_S V
}\tag{xx}
\]
where $\lambda(g,w)=(g,g\cdot w)$, hence $q(g,w)=(\bar g,g\cdot w)$, where
$\bar g$ denotes the image of $g$ under the canonical map
$G(S')\to(G/N)(S')$. Finally, one sees from diagram (x) that $Z\to X$ makes
$Z$ a principal bundle with base $X$ and group $N$ acting on the right by
$(g,v)\cdot n=(gn,v)$, so that in (xx), $q:G\times_S W\to X$ makes
$G\times_S W$ a principal bundle with base $X$ and group $N$ acting on the
right by
\[
 (g,w)\cdot n=(gn,n^{-1}\cdot w).
\]
One summarizes the main preceding morphisms in the following diagram:
\[
\xymatrix@C=2.4em@R=2.2em{
 &G\times_S W\ar[d]^q\ar[drr]^\varphi&&\\
 G/N&X\ar[l]_p\ar[rr]^\psi\ar@{^{(}->}[d]_i&&V\\
 &G/N\times_S V\ar@{-->}[ul]^{\mathrm{pr}_1}\ar@{-->}[urr]_{\mathrm{pr}_2}&&
}\tag{D}
\]
where $\psi=\mathrm{pr}_2\circ i$ and $\varphi=\psi\circ q$, i.e.
$\varphi(g,w)=g\cdot w$. If $v$ is a section of $V$ over $S$, the subsheaf
$X_v$ of $X$ which is the inverse image of this section under $\psi$ is given
by $X_v(S')=$ the set of subsheaves $W'$ of $V_{S'}$ which are locally
conjugate to $W_{S'}$ under $G$, and which contain the section $v_{S'}$ of
$V_{S'}$, while the subsheaf of $G\times_S W$ which is the inverse image of
$v$ under $\varphi$ is isomorphic to the subsheaf $M_v$ of $G$ given by
$M_v(S')=$ the set of $g\in G(S')$ such that $v_{S'}\in g\cdot W(S')$, i.e.
such that $g^{-1}v_{S'}\in W(S')$. If $v$ is a section of $W$ and not merely
of $V$, then $M_v$ clearly contains $N$.

In these descriptions, one has not used the fact that $G,V,W$ were
representable (nor that the site on which one works is defined in terms of
preschemes!). But suppose now that $N$ is representable and faithfully flat
and quasi-compact over $S$, and that $G/N$ is representable. When $S$ is the
spectrum of a field, and $G$ is of finite type over $k$, one knows that this
hypothesis is necessarily satisfied (VIII 6 and VI\textsubscript{B}.11.18).
One then sees from the cartesian diagram (xx), using faithfully flat
quasi-compact descent theory and the fact that $Z\to G\times_S V$ is a closed
immersion, that $X$ is representable (it is obtained by descent of the closed
subprescheme $Z$ of $G\times_S V$ along the faithfully flat and quasi-compact
morphism $G\times_S V\to G/N\times_S V$). Thus diagram (D) is a diagram of
morphisms of preschemes over $S$.

From now on one assumes that $S$ is the spectrum of a field $k$, and that
$G,V,W$ are of finite type over $k$. Let $\mathfrak n$ be the Lie algebra of
$N$, so that $\dim N\leq\operatorname{rank}\mathfrak n$, equality holding if
and only if $N$ is smooth over $k$ (Exp VI\nde{Having been unable to identify
this reference, we refer to Theorem II.5.2.1 of the book:
M. Demazure \& P. Gabriel, \emph{Groupes alg\'ebriques I}, Masson (1970).}).
Let $a\in W(k)$, and consider the subscheme $M_a$ of $G$ defined above,
containing $N$, and isomorphic to $\varphi^{-1}(a)$; one shall denote by
$\mathfrak m_a$ its Zariski tangent space at the identity element $e$ of $N$,
so that one has
\[
 \mathfrak n\subset\mathfrak m_a,\qquad
 \dim N\leq\operatorname{rank}\mathfrak n\leq\mathfrak m_a.\tag{1}
\]
% typo?: kept as printed: the last term of (1) has no rank operator.

\begin{lemma}{1.1}\label{XIII.1.1}
With the preceding notation:

\noindent a) Consider the following conditions:
\begin{enumerate}[label=\textup{(\roman*)}]
\item $\mathfrak n=\mathfrak m_a$ and $N$ is smooth over $k$.
\item[\textup{(i bis)}] $\dim N=\operatorname{rank}_k\mathfrak m_a$.
\item $\psi:X\to V$ is unramified at $(\bar e,a)$.
\item $M_a$ and $N$ coincide in a neighborhood of $e$.
\end{enumerate}
Then one has the implications
$\textup{(i)}\Leftrightarrow\textup{(i bis)}\Rightarrow\textup{(ii)}
\Leftrightarrow\textup{(iii)}$.

\noindent b) Suppose that $\varphi:G\times_S W\to V$ is smooth at $(e,a)$.
Then $M_a$ is smooth over $k$ at $e$, and $\psi$ is smooth at $(\bar e,a)$.
\end{lemma}
\begin{proof}
a) The equivalence of (i) and (i bis) follows at once from relations (1) and
from the fact noted above that $N$ is smooth over $k$ if and only if
$\dim N=\operatorname{rank}_k\mathfrak n$. On the other hand, consider the
inclusion morphism $N\to M_a$; it is well known\nde{See, for example,
EGA IV\textsubscript{4}, Th. 17.11.1 d).} that if $N$ is smooth over $k$ at
$e$ and the tangent map at $e$ is surjective, then $N\to M_a$ is smooth at
$e$, and hence (being an immersion) is an isomorphism at $e$, which shows
that (i) implies (iii). To prove the equivalence of (ii) and (iii), consider,
as above, $X_a=\psi^{-1}(a)$ and use the isomorphism
$M_a\simeq\varphi^{-1}(a)\simeq q^{-1}(X_a)$ to obtain a morphism
$p_a:M_a\to X_a$ which makes $M_a$ a principal homogeneous bundle with group
$N_{X_a}$. Consider the following diagram
\[
\xymatrix{
 M_a\ar[d]_{p_a}&N\ar[l]_{j_a}\ar[d]\\
 X_a\ar[r]^{j'_a}&S=\Spec(k),
}
\]
where $\Spec(k)\to X_n$ is defined by the point $(\bar e,a)$ of $X$, and
$j_a:N\to M_a$ is the canonical immersion. Saying that $\psi$ is unramified
at $(\bar e,a)$ means that $j'_a$ is an open immersion, or again that it
induces an isomorphism $S\isomto\Spec(\OO_{X_a,a})$. Since $p_a$ is flat,
it is equivalent to say that the morphism obtained from the preceding one
by the base change $\Spec(\OO_{M_a,e})\to\Spec(\OO_{X_a,b})$ is an
isomorphism; now this resulting morphism is precisely the morphism
$\Spec(\OO_{N,e})\to\Spec(\OO_{M_a,e})$, which proves the equivalence of
(ii) and (iii).
% typo?: kept as printed: X_n, the direction of j'_a, and the subscripts a/b.

One will moreover note that the proof shows that conditions (ii), (iii)
imply the following condition, apparently stronger than (iii):

\noindent (iii bis) $N$ is an open and closed subscheme of $M$.
% typo?: kept as printed: M rather than M_a.

\noindent b) The first assertion follows from the fact that $M_a$ is
isomorphic to $\varphi^{-1}(a)$, the second from the fact that $q$ is flat
and $q(e,a)=(\bar e,a)$.
\end{proof}

\sgasection{2}{Density theorem and theory of regular points of $G$}
One shall apply the constructions and notation of the preceding section
in the case where $G$ is a smooth connected algebraic group over $k$,
where $V=G$ on which $G$ acts by inner automorphisms, and where $W$ is a
smooth connected algebraic subgroup $H$ of $G$. One shall denote by
$\mathfrak g$ the Lie algebra of $G$, by $\mathfrak h$ that of $H$, by $N$
the normalizer of $H$ in $G$, and by $\mathfrak n$ the Lie algebra of $N$.
If $a\in G(k)$, one shall again denote, as in \No 1, by $M_a$ the inverse
of its transporter into $H$, so that if $a\in H(k)$, one has $N\subset M_a$;
in this case, one denotes by $\mathfrak m_a$ the Zariski tangent space of
$M_a$ at the identity element $e$ of $G$. Note that
\[
 \mathfrak h\subset\mathfrak n\subset\mathfrak m_a\subset\mathfrak g
 \qquad\text{for }a\in H(k).
\]
One shall need the

\begin{lemma}{2.0}\label{XIII.2.0}
For $H=N^0$ to hold, it is necessary and sufficient that
$(\mathfrak g/\mathfrak h)^H=0$ (where the left-hand side denotes the
subspace of invariants under the action of $H$ deduced from the adjoint
representation). When this condition is satisfied, $N$ is smooth and one
has $\dim X=\dim G$. In all cases, $\dim X\leq\dim G$, and this inequality
is an equality if and only if $H$ has finite index in $N$.
\end{lemma}
Indeed, one has seen (II 5.2.3 (i)) that $\mathfrak n$ is equal to the
inverse image of $(\mathfrak g/\mathfrak h)^H$ under the morphism
$\mathfrak g\to\mathfrak g/\mathfrak h$, hence
$(\mathfrak g/\mathfrak h)^H=0$ is equivalent to
$\mathfrak h=\mathfrak n$, which is also equivalent ($H$ being a smooth
connected algebraic subgroup of the algebraic group $N$) to $H=N^0$
(cf. VI.2). This clearly implies that $N$ is smooth. On the other hand,
one has
\[
 \dim X=(\dim G-\dim N)+\dim H=\dim G-(\dim N-\dim H),
\]
so that one has
\[
 \dim X\leq\dim G,
\]
equality being attained if and only if $\dim H=\dim N$, i.e. if and only
if $H$ has finite index in $N$. This is the case in particular if $H=N^0$,
which completes the proof of 2.0.

\begin{theorem}{2.1}\label{XIII.2.1}
Let $G$ be a smooth connected algebraic group over the algebraically closed
field $k$, $H$ a smooth connected algebraic subgroup, $N$ its normalizer,
$\mathfrak g,\mathfrak h,\mathfrak n$ the Lie algebras,
$X=G\times^N H$ the scheme (fibered over $G/N$ with typical fiber $H$)
introduced in \No 1, and $\psi:X\to G$ the canonical morphism (whose image
is also the image of $\varphi:G\times H\to G$ defined by
$\varphi(g,h)=\operatorname{int}(g)h=ghg^{-1}$). The following conditions
are all equivalent:
\begin{enumerate}[label=\textup{(\roman*)}]
\item $H$ contains a Cartan subgroup (XII 1) $C$ of $G$.
\item[\textup{(i bis)}] $H$ has the same reductive rank and the same nilpotent
rank (XII 1) as $G$.
\item $H$ contains a maximal torus $T$ of $G$, and
$(\mathfrak g/\mathfrak h)^T=0$.
\item The set of conjugates of $H$ containing a given maximal torus is
finite and nonempty, and $H$ has finite index in its normalizer.
\item There exists $a\in H(k)$ which is contained in only finitely many
conjugates of $H$ (or merely such that $\psi^{-1}(a)$ has an isolated point),
and $H$ has finite index in its normalizer.
\item[\textup{(iv bis)}] The morphism $\psi:X\to G$ is generically
quasi-finite (i.e. there exists a dense open subset of $X$ on which $\psi$
is quasi-finite), and $H$ has finite index in its normalizer.
\item There exists a dense open subset $U$ of $G$ such that for every
$x\in U(k)$, the set of conjugates of $H$ containing $x$ is finite and
nonempty, i.e. $\psi:X\to G$ is dominant and generically quasi-finite.
\item There exists a dense open subset $U$ of $G$ such that every
$x\in U(k)$ is contained in a conjugate of $H$, i.e. $\psi:X\to G$ is dominant.
\item There exists $a\in H(k)$ such that the subspace of $\mathfrak g/\mathfrak h$
of fixed points of $\ad_{\mathfrak g/\mathfrak h}(a)$ is zero.
\end{enumerate}
Moreover, these conditions imply that $H$ is its own connected normalizer,
i.e. $N$ is smooth and $\dim H=\dim N$, and that $\psi:X\to G$ is
generically \emph{\'etale}.
\end{theorem}
\begin{proof}
By 2.0, one has $\dim X\leq\dim G$, with equality if and only if
$\dim H=\dim N$, i.e. $H$ has finite index in $N$. From the inequality
$\dim X\leq\dim G$ it follows that $\psi$ is dominant if and only if it
is dominant and generically quasi-finite, or again if and only if $\psi$
is generically quasi-finite and $\dim X=\dim G$. Since this last equality
also means, by 2.0, that $H$ has finite index in $N$, one has proved the
equivalence of (vi), (v), (iv bis). The equivalence of (iv) and (iv bis)
is immediate.

The equivalence of (i) and (i bis) is immediate from the definitions, and
is left to the reader. On the other hand, if $H$ contains a Cartan subgroup
$C$ of $G$, it contains the maximal torus $T$ of $C$, which is a maximal
torus of $G$. Since $C$ is the centralizer of $T$, its Lie algebra
$\mathfrak c$ is given by
\[
 \mathfrak c=\mathfrak g^T
\]
(II, 5.2.3 (ii)). Thus, since $H\supset C$, whence
$\mathfrak h\supset\mathfrak n$, it follows that
$\mathfrak h\subset\mathfrak g^T$, which is equivalent, by I, 4.7.3,
to the relation
\[
 (\mathfrak g/\mathfrak h)^T=0.\tag{x}
\]
% typo?: kept as printed: h superset n, and h subset g^T.
Conversely, suppose that $H$ contains the maximal torus $T$ and that the
preceding relation holds, i.e. that one has $\mathfrak h\supset\mathfrak n$;
I say that $H$ contains the centralizer $C$ of $T$ (which will establish
(i) $\Leftrightarrow$ (ii)). This follows from the

\begin{lemma}{2.1.1}\label{XIII.2.1.1}
Let $G$ be a smooth algebraic group over the field $k$, $T$ a subgroup
of multiplicative type of $G$, $C$ its connected centralizer (equal to
the centralizer of $T$ if $G$ is connected and $T$ a torus, (XII 6.6 b)),
and $H$ a smooth subgroup of $G$ containing $T$. For $H$ to contain $C$,
it is necessary and sufficient that its Lie algebra $\mathfrak h$ contain
that $\mathfrak c$ of $C$.
\end{lemma}
Indeed, one knows (XI 2.4) that $\Centr_G(T)$ is smooth over $k$, hence
$C$ is smooth over $k$; likewise $\Centr_H(T)$ is smooth over $k$; now
$\Centr_H(T)=\Centr_G(T)\cap H$ has $\mathfrak h\cap\mathfrak c$ as Lie
algebra, so the hypothesis implies that the smooth subgroup $\Centr_H(T)$
of the smooth group $\Centr_G(T)$ has the same Lie algebra, hence it
contains the connected component $C$ of the latter, so $H$ contains $C$.
\hfill Q.E.D.

Let us prove the equivalence of (i) and (iii), which amounts to proving
that if $H$ contains the maximal torus $T$ of $G$, then the condition
$H\supset C$ (which is also equivalent to (x) above, as one has just seen)
is equivalent to the fact that $H$ has finite index in its normalizer and
that the set of conjugates of $H$ containing $T$ is finite. If $H$ contains
$C$, hence if one has (x), then a fortiori
\[
 (\mathfrak g/\mathfrak h)^H=0,\tag{xx}
\]
now one knows that $\mathfrak n$ is the inverse image of the left-hand
side of the preceding relation under the canonical homomorphism
$\mathfrak g\to\mathfrak g/\mathfrak h$ (II 5.2.3 (i)), so the preceding
relation also means, by 2.0, that $H=N^0$; a fortiori $H$ has finite index
in its normalizer. Consider now the diagram of subgroups
\[
\xymatrix{
 T\ar[r]&N(T)\cap H\ar[r]\ar[d]&H\ar[d]\\
 &N(T)\cap N(H)\ar[r]\ar[d]&N(H)\\
 &N(T).
}
\]
Using the conjugacy theorem for maximal tori in $H$ (XII 6.6 a)), one
sees that every conjugate of $H$ containing $T$ is conjugate to $H$ by
an element of $N(T)(k)$, hence that the set of conjugates of $H$ containing
$T$ is in one-to-one correspondence with the set of points of
$N(T)/N(T)\cap N(H)$ with coefficients in $k$; now, since $H\supset C$,
one has $N(T)\cap H\supset C$, so the preceding set is a quotient of
$(N(T)/C)(k)$, which is a finite set, and is therefore finite. This proves
that (i) $\Rightarrow$ (iii). Conversely, suppose (iii), i.e.
$N(T)/N(T)\cap N(H)$ is finite and $N(H)/H$ is finite. Using the conjugacy
theorem in $H$ again, one again sees that the homomorphism
\[
 N(T)\cap N(H)/N(T)\cap H\longrightarrow N(H)/H
\]
induced by the preceding diagram is bijective on the points with values
in $k$ (in fact, it is an isomorphism), so, since the second is finite,
so is the first; hence $N(T)\cap H$ has finite index in $N(T)$, and thus
contains $C=N(T)^0$, so $H\supset C$. Thus (i), (i bis), (ii), (iii)
are equivalent conditions.

Let us prove that (ii) $\Rightarrow$ (vii). One sees at once that conditions
(ii) and (vii) are each invariant under an extension $k\to k'$ of the
base field, with $k'$ algebraically closed, which allows one to suppose
that $k$ has infinite transcendence degree over its prime subfield.
Then it is well known (and one verifies immediately) that there exists
an element $a$ of $T(K)$ such that the subgroup of $T(K)$ which it
generates is dense in $T$ for the Zariski topology.
% typo?: kept as printed: T(K), although the base field is k.
